\documentclass{umk-matematika}

\begin{document}

\maketitlepage
    {Практическая работа №9}
    {Уравнение касательной}
    {}

\section{Фундамент (1--4 балла)}

\textbf{Задача 1 (1 балл).} Найдите значение производной в точке $x_0$: $f(x) = x^2,\ x_0 = 3$

\textbf{Задача 2 (2 балла).} Найдите $f(x_0)$: $f(x) = x^2,\ x_0 = 3$

\textbf{Задача 3 (3 балла).} Запишите общий вид уравнения касательной к графику функции $f(x)$ в точке $x_0$.

\textbf{Задача 4 (4 балла).} Составьте уравнение касательной: $f(x) = x^2,\ x_0 = 3$

\section{Стены (5--7 баллов)}

\textbf{Задача 5 (5 баллов).} Составьте уравнение касательной: $f(x) = x^2 - 4x,\ x_0 = 1$

\textbf{Задача 6 (6 баллов).} Составьте уравнение касательной: $f(x) = x^3,\ x_0 = -1$

\textbf{Задача 7 (7 баллов).} Составьте уравнение касательной: $f(x) = 2x^2 - 3x + 1,\ x_0 = 2$

\section{Кровля (8--10 баллов)}

\textbf{Задача 8 (8 баллов).} Найдите точку на графике $f(x) = x^2$, в которой касательная параллельна прямой $y = 4x - 1$.

\textbf{Задача 9 (9 баллов).} Составьте уравнение касательной к графику $f(x) = x^2 + 2x$ в точке с абсциссой $x_0 = -3$.

\textbf{Задача 10 (10 баллов).} Уклон дорожного покрытия описывается функцией $f(x) = -0{,}5x^2 + 3x$, где x — расстояние от начала участка (м). Найдите уравнение касательной в точке $x_0 = 2$ (это уклон в данной точке).

\newpage
\section*{Ответы}

\begin{enumerate}
  \item $f'(3) = 6$
  \item $f(3) = 9$
  \item $y = f(x_0) + f'(x_0)(x - x_0)$.
  \item $y = 6x - 9$
  \item $y = -2x - 1$
  \item $y = 3x + 2$
  \item $y = 5x - 7$
  \item точка $(2;\ 4)$.
  \item $y = -4x - 9$
  \item $y = x + 2$
\end{enumerate}

\end{document}